\section{Group or moment? The same flight before, during and after}\label{sec:prepost}

Do pilots fly together in the moments when the progress is straight anyway? If so, the same flight
would already be straight just before the group forms, and still just after it breaks.

\paragraph{Windows.} Every stretch of group time is an \emph{event}, with up to four windows of length
$W=\PPWindowMin$~min on the clock of the same flight: \emph{start} and \emph{end}, the first and last
$W$ of the group, when it lasts at least $W$; \emph{before} and \emph{after}, the last $W$ of the solo
stretch that precedes the group and the first $W$ of the one that follows it, when that solo stretch
lasts at least $W$, whatever the length of the group. A window of 50~min holds about 200 origins at
$\tau=1000$~s. For two windows the ensemble MSD is taken on the events that have both, and their
difference is read on $\Hloc$ and on the phase statistics. Entering is start minus before; leaving is
end minus after, group minus alone in both cases. The two readings predict different things:
\begin{itemize}
  \item if the straighter progress belongs to the group, it appears on entering and on leaving, and
        before and after are alike;
  \item if it belongs to the moment, before is already high, and the changes are those of any two
        consecutive windows.
\end{itemize}

\paragraph{Controls.} Two consecutive windows with no change of state, drawn in the solo time of the
group flights ($\A$) and in solo flights ($\Sl$), measure how much a flight drifts from one window to
the next, later minus earlier. Entering has the same time order as the controls; leaving has the
opposite one, so on leaving the drift enters with the opposite sign.

\begin{figure}
\centering
\includegraphics[width=\textwidth]{prepost_3000_curves.pdf}
\caption{$\Hloc(\tau)$ of the same flights alone (green) and in the group (orange), entering (left)
and leaving (middle), 5--95\% bands. Right: the differences, group minus alone, with the two controls
(later minus earlier window; dashed); for leaving, the drift to compare is the control with the
opposite sign.}
\label{fig:prepost}
\end{figure}

\begin{figure}
\centering
\includegraphics[width=0.86\textwidth]{prepost_3000_profile.pdf}
\caption{The four windows on the events that have all of them. Start and end overlap when the group
lasts less than $2W$.}
\label{fig:prepost-profile}
\end{figure}

\begin{table}
\centering\footnotesize\setlength{\tabcolsep}{4pt}
\begin{tabular}{@{}l r c c c@{}}
\toprule
comparison & events & $\Delta H$ & $\Delta H_\mathrm{loc}(100)$ & $\Delta H_\mathrm{loc}(1000)$ \\
\midrule
entering: start $-$ before & 6\,524 & \makecell{$+$0.010\\{\scriptsize[$+$0.009, $+$0.011]}} & \makecell{$+$0.005\\{\scriptsize[$+$0.004, $+$0.006]}} & \makecell{$+$0.027\\{\scriptsize[$+$0.023, $+$0.031]}} \\[2pt]
leaving: end $-$ after & 13\,643 & \makecell{$+$0.005\\{\scriptsize[$+$0.004, $+$0.006]}} & \makecell{$-$0.001\\{\scriptsize[$-$0.002, $-$0.000]}} & \makecell{$+$0.031\\{\scriptsize[$+$0.028, $+$0.033]}} \\[2pt]
after $-$ before & 6\,320 & \makecell{$+$0.004\\{\scriptsize[$+$0.003, $+$0.005]}} & \makecell{$+$0.006\\{\scriptsize[$+$0.005, $+$0.007]}} & \makecell{$-$0.003\\{\scriptsize[$-$0.007, $+$0.001]}} \\[2pt]
control A: later $-$ earlier & 16\,998 & \makecell{$+$0.005\\{\scriptsize[$+$0.004, $+$0.005]}} & \makecell{$+$0.005\\{\scriptsize[$+$0.005, $+$0.006]}} & \makecell{$-$0.002\\{\scriptsize[$-$0.005, $+$0.000]}} \\[2pt]
control S: later $-$ earlier & 69\,648 & \makecell{$+$0.009\\{\scriptsize[$+$0.008, $+$0.009]}} & \makecell{$+$0.010\\{\scriptsize[$+$0.010, $+$0.011]}} & \makecell{$-$0.002\\{\scriptsize[$-$0.003, $-$0.000]}} \\[2pt]
\bottomrule
\end{tabular}

\caption{Differences of the fitted $H$ and of $\Hloc$ at 100 and 1000~s (value, 5--95\% band) between
windows, on the events with both windows. Controls: two consecutive windows with no change of state,
later minus earlier, in the solo time of the group flights ($\A$) and in solo flights ($\Sl$).}
\label{tab:prepost}
\end{table}

\paragraph{Results at 1000~s.} The first prediction holds. $\Delta\Hloc(1000~\mathrm{s})$ is
\PPEntryDHloc{} \PPEntryDHlocBand{} entering and \PPExitDHloc{} \PPExitDHlocBand{} leaving, against
drifts of \PPCtrlADHloc{} and \PPCtrlSDHloc{} in the controls, and after minus before is
\PPAfterBeforeDHloc{} \PPAfterBeforeDHlocBand. On the \PPProfileEvents{} events with all four
windows, $\Hloc(1000~\mathrm{s})$ is \PPProfileHlocPre{} before, \PPProfileHlocStart{} and
\PPProfileHlocEnd{} in the group and \PPProfileHlocPost{} after (Figure~\ref{fig:prepost-profile}):
a step close to the within-flight estimate of Section~\ref{sec:other} (\DHlocThousandFlight).

\paragraph{Results at short lags and on the phases.} Entering, $\Hloc(100~\mathrm{s})$ and the fitted
$H$ change by \PPEntryDHlocHundred{} and \PPEntryDH, about the drift of the controls
(\PPCtrlADHlocHundred{} and \PPCtrlSDHlocHundred; \PPCtrlADH{} and \PPCtrlSDH). Leaving, they change by
\PPExitDHlocHundred{} and \PPExitDH, which exceed the drift taken with the opposite sign: a small
asymmetry at short lags, which we do not interpret. Glides are faster and flatter in the group in
both directions: glide speed \PPEntrySpeed{} km/h entering and \PPExitSpeed{} leaving, against drifts
of \PPCtrlASpeed{} and \PPCtrlSSpeed{} km/h; height lost per km \PPEntrySink{} entering and
\PPExitSink{} leaving, against drifts of \PPCtrlASink{} and \PPCtrlSSink{} m/km, so that once each
drift is taken with its sign the glide flattens by similar amounts both ways (Table~\ref{tab:prepost-phase},
Appendix~\ref{app:tables}). The climb statistics rise on entering beyond the drift and do not fall on
leaving beyond it.

\paragraph{Limits.} The windows fix the pilot, the glider, the day and the slow drift of the flight,
not what brings pilots together: a convergence line, a ridge or a common route may last as long as
the group and straighten the progress by itself, and two pilots must move alike to stay within
500~m for long. The events are long groups with solo flight beside them (\PPEntryEvents{} entries,
\PPExitEvents{} exits). With windows of \PPShortWindowMin~min the step at 1000~s is not seen
(\PPShortEntryDHloc{} entering, \PPShortExitDHloc{} leaving, against drifts of \PPShortCtrlADHloc{} and
\PPShortCtrlSDHloc): such a window holds fewer than thirty origins at 1000~s, so the test at that lag
needs windows of \PPWindowMin~min (Appendix~\ref{app:prepost}).

\paragraph{What we keep.} With windows of \PPWindowMin~min, on the same flight the straighter
progress over a quarter of an hour starts with the group and stops with it, and the glides change
with it both ways; the moment of joining does not explain them. The climbs change when the group
forms and not when it breaks, so for them the moment remains a possible cause.
